% T12 C/D, chapitre 6 : Applications affines et isométries. Source : le livre « Terminales scientifiques » (Affine/Lesona, Isometrie/Lesona, Isometrie/Exo).
% Les exercices tirés d’autres ouvrages ne sont pas repris (décision du 2 novembre 2026).

\TCDLessonStyleOn
\providecommand{\Plan}{\mathscr{P}}
\providecommand{\Espace}{\mathcal{E}}
\providecommand{\bary}[1]{\mathrm{Bar}\left\lbrace #1 \right\rbrace}
\providecommand{\Id}{\mathrm{Id}}
\chapter{Applications affine}
      Dans ce chapitre, $\Plan$ d\'esigne le plan affine et $\Espace$ l'espace affine.
      \section{D\'efinitions et exemples}
            \begin{defi}
\index{Application affine}
                  Une application $f : \Plan \rightarrow \Plan$ est une \emph{application affine}
                  si elle conserve le barycentre.
                  
                  Si $G = \bary{(A,\alpha) ; (B,\beta)}$.
                  En posant $G' = f(G)$, $A' = f(A)$ et $B' = f(B)$, on a 
                  $G' = \bary{(A',\alpha) ; (B',\beta)}$.
            \end{defi}
            %
            \begin{ex}
                  \begin{enumerate}
                        \item Consid\'erons la translation $t_{\vec{u}}$ de vecteur $\vec{u}$.
                              En posant $t_{\vec{u}} (M) = M'$, on a $\overrightarrow{MM'} = \vec{u}$.
                              
                              Soit $G = \bary{(A,\alpha) ; (B,\beta)}$.
                              Soient
                              $G' = t_{\vec{u}} (G)$,
                              $A' = t_{\vec{u}} (A)$ et
                              $B' = t_{\vec{u}} (B)$
                              .
                              On a alors
                              \begin{align*}
                              \alpha \overrightarrow{G'A'} + \beta \overrightarrow{G'B'}
                                    & =  \alpha (\overrightarrow{G'G} + \overrightarrow{GA} + \overrightarrow{AA'})
                                        +\beta (\overrightarrow{G'G} + \overrightarrow{GB} + \overrightarrow{BB'})      \\
                                    & =  \alpha (-\vec{u} + \overrightarrow{GA} + \vec{u})
                                        +\beta (-\vec{u} + \overrightarrow{GB} + \vec{u})      \\
                                    & = \alpha \overrightarrow{GA} + \beta \overrightarrow{GB}      \\
                                    & = \vec{0}.
                              \end{align*}
                              Donc $G' = \bary{(A';\alpha);(B';\beta)}$ et ainsi, $t_{\vec{u}}$ est une application affine.
                  \end{enumerate}
            \end{ex}
            %
            \begin{ex}
                  Consid\'erons l'homoth\'etie $h$ de centre $I$ et de rapport $k$.
                  En posant $h(M) = M'$, on a $\overrightarrow{IM'} = k \overrightarrow{IM}$.
                  
                  On veut montrer que $h$ est une application affine.
                  Consid\'erons le point
                  \[
                  G = \bary{(A,\alpha);(B,\beta)}.
                  \]
                  Posons $G' = h(G)$, $A' = h(A)$ et $B' = h(B)$.
                  On a alors 
                  \begin{align*}
                        \alpha \overrightarrow{G'A'} + \beta \overrightarrow{G'B'}
                              & =  \alpha (\overrightarrow{G'I} + \overrightarrow{IA'})
                                  +\beta (\overrightarrow{G'I} + \overrightarrow{IB'})      \\
                              & = \alpha (k \overrightarrow{GI} + k \overrightarrow{IA}) 
                                  +\beta (k \overrightarrow{GI} + k \overrightarrow{IB})      \\
                              & = k \left[ \alpha \overrightarrow{GA} + \beta \overrightarrow{GB} \right]       
                                = k \vec{0}      
                                = \vec{0}.
                  \end{align*}
            \end{ex}
            %
            \begin{ex}
                  Consid\'erons la r\'eflexion d'axe $(\mathscr{D})$.
                  On veut montrer que la r\'eflexion est une application affine.
                  
                        \begin{tikzpicture}[line cap=round,line join=round,>=triangle 45,x=2.0cm,y=2.0cm]
                              \clip(-0.94,-0.73) rectangle (2.88,2.42);
                              \draw[fill=black,fill opacity=0.1] 
                                    (0.8,1.4) -- (0.81,1.19) -- (1.03,1.19) -- (1.02,1.41) -- cycle; 
                              \draw[fill=black,fill opacity=0.1] 
                                    (0.82,0.7) -- (0.83,0.49) -- (1.05,0.49) -- (1.04,0.71) -- cycle; 
                              \draw[fill=black,fill opacity=0.1] 
                                    (0.85,-0.14) -- (0.85,-0.36) -- (1.07,-0.35) -- (1.06,-0.14) -- cycle; 
                              \draw [domain=-0.94:2.88] plot(\x,{(--1.91-1.8*\x)/0.05});
                              \draw (0.42,-0.15)-- (0,0.68);
                              \draw [dash pattern=on 1pt off 1pt] (0,0.68)-- (-0.36,1.37);
                              \draw (1.7,-0.12)-- (2.08,0.74);
                              \draw [dash pattern=on 1pt off 1pt] (2.08,0.74)-- (2.4,1.45);
                              \draw [dotted] (-0.36,1.37)-- (1.02,1.41);
                              \draw [] (0.33,1.44) -- (0.33,1.35);
                              \draw [dotted] (1.02,1.41)-- (2.4,1.45);
                              \draw [] (1.71,1.48) -- (1.71,1.38);
                              \draw [dotted] (2.08,0.74)-- (1.04,0.71);
                              \draw [] (1.58,0.68) -- (1.58,0.77);
                              \draw [] (1.54,0.68) -- (1.54,0.77);
                              \draw [dotted] (0,0.68)-- (1.04,0.71);
                              \draw [] (0.5,0.74) -- (0.5,0.65);
                              \draw [] (0.53,0.74) -- (0.54,0.65);
                              \draw [dotted] (1.06,-0.14)-- (0.42,-0.15);
                              \draw [] (0.78,-0.19) -- (0.78,-0.1);
                              \draw [] (0.75,-0.19) -- (0.74,-0.1);
                              \draw [] (0.71,-0.19) -- (0.71,-0.1);
                              \draw [dotted] (1.7,-0.12)-- (1.06,-0.14);
                              \draw [] (1.42,-0.17) -- (1.42,-0.08);
                              \draw [] (1.39,-0.17) -- (1.38,-0.08);
                              \draw [] (1.35,-0.17) -- (1.35,-0.08);
                              \draw (1.2,2.24) node[anchor=north west] {$ (\mathscr{D}) $};
                              \begin{scriptsize}
                              \fill [color=black] (0.42,-0.15) circle (1.5pt) node[below left] {$A$};
                              \fill [color=black] (0,0.68) circle (1.5pt) node[below left] {$B$};
                              \fill [color=black] (-0.36,1.37) circle (1.5pt) node[below left] {$G$};
                              \fill [color=black] (1.7,-0.12) circle (1.5pt) node[below right] {$A'$};
                              \fill [color=black] (2.08,0.74) circle (1.5pt) node[below right] {$B'$};
                              \fill [color=black] (2.4,1.45) circle (1.5pt) node[below right] {$G'$};
%                              \fill [color=black] (1.02,1.41) circle (1.5pt) node {$C$};
%                              \fill [color=black] (1.04,0.71) circle (1.5pt) node {$D$};
%                              \fill [color=black] (1.06,-0.14) circle (1.5pt) node {$E$};
                              \end{scriptsize}
                        \end{tikzpicture}
                  
                  Soit $G = \bary{(A,\alpha);(B,\beta)}$.
                  Alors, $A$, $G$ et $B$ sont align\'es.
                  \begin{align*}
                        \alpha \overrightarrow{GA} + \beta \overrightarrow{GB} & = \vec{0} \\
                        \alpha \overrightarrow{GA} + \beta \overrightarrow{GA} + \beta \overrightarrow{AB} & = \vec{0} \\
                        (\alpha+\beta) \overrightarrow{GA} + \beta \overrightarrow{AB} & = \vec{0}
                  \end{align*}
                  et donc $\overrightarrow{AG} = \frac{\beta}{\alpha + \beta} \overrightarrow{AB}$.
                  Comme
                  $\overline{AG} = \frac{\beta}{\alpha + \beta} \overline{AB}$, on a 
                  $\frac{\overline{AG}}{\overline{AB}} = \frac{\beta}{\alpha + \beta}$.
                  
                  De plus, $(AA') \parallel (BB') \parallel (GG')$,
                  donc 
                  $\frac{\overline{AG}}{\overline{AB}} = \frac{\overline{A'G'}}{\overline{A'B'}}$.
                  D'o\`u
                  $\frac{\overline{A'G'}}{\overline{A'B'}} = \frac{\beta}{\alpha+\beta}$
                  et, par cons\'equent,
                  $\overrightarrow{A'G'} = \frac{\beta}{\alpha+\beta} \overrightarrow{A'B'}$.
                  
                  Ainsi, $G' = \bary{(A',\alpha);(B',\beta)}$. \qed
            \end{ex}
            %
            \begin{ex}
                  Consid\'erons la rotation $r$ de centre $I$ et d'angle $\alpha$.
                  Elle a alors pour expression complexe $z' = az + b$ o\`u $|a| = 1$.
                  
                  Soit $G = \bary{(A,\alpha);(B,\beta)}$. 
                  Posons
                        $$
                        G' = r(G)
                              \qquad
                        A' = r(A)
                              \qquad
                        B' = r(B).
                        $$
                  L'affixe de $\alpha \overrightarrow{G'A'} + \beta \overrightarrow{G'B'}$ 
                  est $\alpha (z_A' - z_G') + \beta (z_B' - z_G')$,
                  ce qui est \'egal \`a 
                  \[
                        \alpha (a z_A + b - a z_G - b) + \beta (a z_B + b - a z_G - b) =
                                                      a \left[ \alpha (z_A - z_G) + \beta (z_B - z_G) \right] \]
                  qui n'est autre que l'affixe de 
                  $a ( \underbrace{\alpha \overrightarrow{GA} + \beta \overrightarrow{GB}}_{\vec{0}} )$.
                  %
                  Ainsi, $G'= \bary{(A,\alpha);(B,\beta)}$.
            \end{ex}
%
      \begin{thm}
            La composition de deux applications affines est une application affine.
            \newline \smallnota{
            C'est-\`a-dire que si $f$ et $g$ sont deux applications affines, alors $f\circ g$ est aussi une application affine.
            }
      \end{thm}
%
      \subsection{Application lin\'eaire associ\'ee}
\index{Application affine!Application lineaire associee@Application lin\'eaire associ\'ee}
            Soit $f : \Plan \rightarrow \Plan , M \mapsto M'$ une application affine.
            L'application $\Phi : \overrightarrow{AB} \mapsto \Phi(\overrightarrow{AB}) = \overrightarrow{A'B'}$
            est appel\'ee \emph{application lin\'eaire associ\'ee \`a l'application affine $f$}.
            %
            \begin{ex}
                  Trouvons l'application lin\'eaire associ\'ee \`a une translation $t_{\vec{u}}$.
                  
                        \begin{tikzpicture}[line cap=round,line join=round,>=triangle 45,x=1.5cm,y=1.5cm]
%                              \clip(-2.28,-0.3) rectangle (3.14,1.43);
                              \draw [->] (-0.37,0.19) -- (0.91,0.09);
                              \draw [->] (-1.91,0.12) -- (-0.36,0.95);
                              \draw [->] (-0.37,0.19) -- (0.91,0.09);
                              \draw [->] (1.19,1.03) -- (2.47,0.92);
                              \begin{scriptsize}
                              \fill [color=black] (-0.37,0.19) circle (1.5pt);
                              \draw[color=black] (-0.29,0.32) node {$A$};
                              \fill [color=black] (0.91,0.09) circle (1.5pt);
                              \draw[color=black] (0.99,0.22) node {$B$};
                              \draw[color=black] (-0.84,0.65) node[below] {$\vec{u}$};
                              \fill [color=black] (1.19,1.03) circle (1.5pt);
                              \draw[color=black] (1.28,1.16) node {$A'$};
                              \fill [color=black] (2.47,0.92) circle (1.5pt);
                              \draw[color=black] (2.56,1.06) node {$B'$};
                              \end{scriptsize}
                        \end{tikzpicture}

                  Appelons $T$ l'application lin\'eaire associ\'ee \`a $t_{\vec{u}}$.
                  On a alors $T : \overrightarrow{AB} \mapsto \overrightarrow{A'B'} = \overrightarrow{AB}$.
                  Ainsi, $T = \Id$.
            \end{ex}
            %
            \begin{ex}
                  Cherchons l'application lin\'eaire associ\'ee \`a une homoth\'etie $h(I,k)$.
                  
                        \begin{tikzpicture}[line cap=round,line join=round,>=triangle 45,x=2.0cm,y=2.0cm]
                              \clip(-1.94,-0.83) rectangle (1.37,1.12);
                              \draw [->] (-1.28,0.05) -- (-0.4,0.8);
                              \draw [->] (-0.8,-0.72) -- (1.15,0.94);
                              \draw [dotted] (-1.69,0.69)-- (1.15,0.94);
                              \draw [dotted] (-1.69,0.69)-- (-0.8,-0.72);
                              \begin{scriptsize}
                              \fill [color=black] (-1.28,0.05) circle (1.5pt) node[above] {$A$};
                              \fill [color=black] (-0.4,0.8) circle (1.5pt) node[above] {$B$};
                              \fill [color=black] (-1.69,0.69) circle (1.5pt) node[above] {$I$};
                              \fill [color=black] (-0.8,-0.72) circle (1.5pt) node[above] {$A'$};
                              \fill [color=black] (1.15,0.94) circle (1.5pt) node[above] {$B'$};
                              \end{scriptsize}
                        \end{tikzpicture}
                                          
                  Soit $H$ l'application lin\'eaire associ\'ee \`a $h(I,k)$.
                  On a $H : \overrightarrow{AB} \mapsto \overrightarrow{A'B'} = k \overrightarrow{AB}$.
                  
                  L'application $H$ est appel\'ee une \emph{homoth\'etie vectorielle}.
\index{Homothetie vectorielle@Homoth\'etie vectorielle}
            \end{ex}
%
      \subsection{Expression analytique d'une application affine}
\index{Expression analytique}
            \begin{defi}[Expression analytique]
                  Soit $f$ l'application affine qui \`a tout point $M(x,y)$ associe le point $M'(x',y')$.
                  On appelle \emph{expression analytique} de $f$ la donn\'ee de $(x',y')$ en fonction de $(x,y)$.
                  Elle est souvent donn\'ee sous la forme : 
                  \begin{equation}
                        \left\lbrace
                              \begin{array}{l}
                                    x' = a_1 x + a_2 y + a_3 \\
                                    y' = b_1 x + b_2 y + b_3 
                              \end{array}
                        \right.
                        \label{eq:expression analytique}
                  \end{equation}
                  o\`u $a_1$, $a_2$, $a_3$, $b_1$, $b_2$ et $b_3$ sont des r\'eels.
            \end{defi}
            %
            \begin{thm}
                  Une application affine $f$ du plan $\Plan$ est enti\`erement d\'efinie par les donn\'ees 
                  des transform\'es de trois points non-align\'es $A$, $B$ et $C$.
			      \begin{proof} ~\par
			            \TCDSideBySide[.64]{%
			                  Soient $A$, $B$ et $C$ trois points non-align\'es du plan.
			                  Alors $(A ; \overrightarrow{AB} , \overrightarrow{AC})$ forme un rep\`ere.
      
                              Pour tout point $M$ du plan, il existe $(x,y)$ tel que 
                              $\overrightarrow{AM} = x \overrightarrow{AB} + y \overrightarrow{AC}$.
			                  Comme $f$ est une application affine,
			                  $\overrightarrow{A'M'} = x \overrightarrow{A'B'} + y \overrightarrow{A'C'}$,
			                  et donc
			            }{%
			                  \resizebox{\linewidth}{!}{%
			                  \begin{tikzpicture}[>=triangle 45]
                                    \draw[->] (0,0) -- (3,-1);
                                    \draw[->] (0,0) -- (2,1);
                                    \draw[fill=loko2] (0,0) circle (1pt) node[left] {$A$};
                                    \draw[fill=loko2] (3,-1) circle (1pt) node[right] {$B$};
                                    \draw[fill=loko2] (2,1) circle (1pt) node[right] {$C$};
			                  \end{tikzpicture}%
			                  }
			            }
                        \[
                        \begin{pmatrix} x' - x_A' \\ y' - y_A' \end{pmatrix}
                        =
                        x \begin{pmatrix} x_B' - x_A' \\ y_B' - y_A' \end{pmatrix}
                        +
                        y \begin{pmatrix} x_C' - x_A' \\ y_C' - y_A' \end{pmatrix}.
                        \]
%                        Il vient que
%                        \[
%                        \left\lbrace
%                        \begin{array}{l}
%                        x' = x (x_B'-x_A')+y(x_C'-x_A')+x_A' \\
%                        y' = y (y_B'-y_A')+y(y_C'-y_A')+y_A'
%                        \end{array}
%                        \right.
%                        \]
                        Ce qui d\'etermine $(x',y')$.
                  \end{proof}
            \end{thm}
            %
%            \begin{exe}
%                  On rapporte le plan $\Plan$ \`a un rep\`ere orthonorm\'e $(O;\ivec,\jvec)$.
%                  On donne les points $A(1,2)$, $B(3,0)$ et $C(-1,0)$.
%                  Soit $f$ l'application affine telle que $f(A) = A'(3,4)$, $f(B) = B'(5,2)$ et $f(C) = A$.
%                  \begin{enumerate}
%                        \item Ecrire l'expression analytiques de $f$ 
%                              %dans le rep\`ere $(A;\overrightarrow{AB},\overrightarrow{AC})$
%                              .
%                        \item On consid\`ere $(\mathscr{D}) : y = x + 1$ et $(\Delta) : y = - x + 1$.
%                              D\'eterminer les tranform\'es de $(\mathscr{D})$ et de $(\Delta)$ par $f$.
%                  \end{enumerate}
%                  %
%                  \begin{valiny}
%                        \begin{tikzpicture}[line cap=round,line join=round,>=stealth,x=1.0cm,y=1.0cm]
%                        \clip (-3.5,-1.5) rectangle (5,4.5);
%                        \draw[->,color=black] (-5.25,0) -- (4.74,0);
%                        \foreach \x in {-5,-4,-3,-2,-1,1,2,3,4}
%                        \draw[shift={(\x,0)},color=black] (0pt,2pt) -- (0pt,-2pt) node[below] {\footnotesize $\x$};
%                        \draw[->,color=black] (0,-4.04) -- (0,4.19);
%                        \foreach \y in {-4,-3,-2,-1,1,2,3,4}
%                        \draw[shift={(0,\y)},color=black] (2pt,0pt) -- (-2pt,0pt) node[left] {\footnotesize $\y$};
%                        \draw[color=black] (0pt,-10pt) node[right] {\footnotesize $0$};
%                        \clip(-5.25,-4.04) rectangle (10.74,4.19);
%                        \draw [domain=-5.25:10.74] plot(\x,{(--6-2*\x)/2});
%                        \draw [->,thick] (1,2) -- (-1,0);
%                        \draw [->,thick] (1,2) -- (3,0);
%                        \draw [->,thick] (0,0) -- (0,1);
%                        \draw [->,thick] (0,0) -- (1,0);
%                        \draw [domain=-5.25:10.74,loko1] plot(\x,{(-2-2*\x)/-2});
%                        \draw [loko1] (2.29,3.21) node[anchor=north west] {$ (\mathscr{D}) $};
%                        \draw [domain=-5.25:10.74,loko2] plot(\x,{(--1-1*\x)/1});
%                        \draw [loko2] (-2.5,3) node[anchor=north west] {$ (\Delta) $};
%                        \begin{scriptsize}
%                        \fill [color=black] (1,2) circle (1.5pt) node[above] {$A$};
%                        \fill [color=black] (3,0) circle (1.5pt) node[above right] {$B$};
%                        \fill [color=black] (-1,0) circle (1.5pt) node[above left] {$C$};
%                        \draw[color=black] (0,0.5) node[right] {$\vec{\jmath}$};
%                        \draw[color=black] (0.5,0) node[below] {$\vec{\imath}$};
%                        \end{scriptsize}
%                        \end{tikzpicture}
%                        %
%                        \begin{enumerate}
%                              \item Posons $M(x,y)$ et $M'(x',y') = f(M)$.
%                                    On a $\overrightarrow{A'M'} = x \overrightarrow{A'B'} + y \overrightarrow{A'C'}$.
%                                    L'expression analytique de $f$ dans 
%                                    $(A;\overrightarrow{AB},\overrightarrow{AC})$ est alors donn\'ee par 
%                                    \[
%                                    \left\lbrace
%                                          \begin{array}{l}
%                                                x'=x(x_B'-x_A')+y(x_C'-x_B')+x_A' \\
%                                                y'=x(y_B'-y_A')+y(y_C'-y_B')+y_A' .
%                                          \end{array}
%                                    \right.
%                                    \]
%                                    Ainsi, l'expression analytique de $f$ %dans $(A;\overrightarrow{AB},\overrightarrow{AC})$ 
%                                    est donn\'ee par :
%                                    \mybox{
%                                    \[
%                                    f :
%                                    \left\lbrace
%                                          \begin{array}{l}
%                                                x' = 2 x - 2 y + 3  \\
%                                                y' = -2 x - 2 y + 4 .
%                                          \end{array}
%                                    \right.
%                                    \]
%                                    }
%                              \item $(\mathscr{D}) : x = 0$ dans $(A ; \overrightarrow{AB} , \overrightarrow{AC})$.
%                                    Donc 
%                                    $
%                                    \left\lbrace
%                                          \begin{array}{l}
%                                                x' = -2y + 3 \\
%                                                y' = -2y + 4 .
%                                          \end{array}
%                                    \right.
%                                    $
%                                    .
%                                    Ainsi, $x'-y' = -1$ et donc $(\mathscr{D'}) : y = x + 1$.
%                                    
%                                    Dans le rep\`ere $(A;\overrightarrow{AB},\overrightarrow{AC})$, 
%                                    $(\Delta) : y = \frac{1}{2}$.
%                                    Ainsi, 
%                                    $
%                                    \left\lbrace
%                                          \begin{array}{l}
%                                                x' = 2x+2 \\
%                                                y' = -2x+3 .
%                                          \end{array}
%                                    \right.
%                                    $
%                                    ,
%                                    et donc $x'+y' = 5$.
%                                    La droite $(\Delta')$ \`a donc pour \'equation $y'=-x+5$ 
%                                    dans le rep\`ere $(A;\overrightarrow{AB},\overrightarrow{AC})$.
%                        \end{enumerate}
%                  \end{valiny}
%            \end{exe}























\providecommand{\Id}{\mathrm{Id}}
\providecommand{\Droite}{(\mathscr{D})}
\providecommand{\DroiteU}{(\mathscr{D}_1)}
\providecommand{\DroiteD}{(\mathscr{D}_2)}
\providecommand{\Arg}[1]{\mathrm{Arg}\left(#1\right)}

\section{Isom\'etrie affine}
      \begin{defi}[Isom\'etrie]
\index{Isometrie@Isom\'etrie}
            Une transformation $f$ est une isom\'etrie si elle conserve la distance, autrement dit :
            si 
            $ \begin{multiform}
                  f(M) = M' \\
                  f(N) = N'
            \end{multiform}      
            $
            alors 
            $MN = MN'$.
            
            Nous admettons que toutes \emph{isom\'etrie} est une \emph{application affine bijective}.
      \end{defi}
      %
      \begin{ex}
            \begin{description}
                  \item[Translation :]
\index{Isometrie@Isom\'etrie!translation}
                        Soit $\vec{u}$ un vecteur et soient $M$ et $N$ deux points du plan.
                        Notons par $t_{\vec{u}}$ la translation de vecteur $\vec{u}$. 
                        On a alors 
                        \begin{IEEEeqnarray*}{rCl}
                              t_{\vec{u}} (M) = M' & \Rightarrow & \overrightarrow{MM'} = \vec{u} \\
                              t_{\vec{u}} (N) = N' & \Rightarrow & \overrightarrow{NN'} = \vec{u} \\
                        \end{IEEEeqnarray*}
                        Ainsi,
                        \[
                              \overrightarrow{M'N'} = \underbrace{\overrightarrow{M'M}}_{-\vec{u}}
                                                      + \overrightarrow{MN}
                                                      + \underbrace{\overrightarrow{NN'}}_{\vec{u}}
                                       \Rightarrow
                                                      \overrightarrow{M'N'} = \overrightarrow{MN}
                        \]
                        et donc $M'N' = MN$.
                  \item[Rotation :]
\index{Isometrie@Isom\'etrie!rotation}
                        Consid\'erons une rotation $r$ de centre $I$.
                        Si $M' = r(M)$ et $N' = r(M)$ alors les triangles $IMN$ et $IM'N'$ sont superposable.
                        Ainsi, $MN = M'N'$.
                        \begin{figure}[H]
                        \centering
                              \begin{tikzpicture}[line cap=round,line join=round,>=triangle 45,x=1.0cm,y=1.0cm]
                              \clip(-2.12,-3.68) rectangle (4.82,2.36);
                              \draw [shift={(1.03,0.24)},color=loko2!80!black,fill=loko2!80!black,fill opacity=0.1] 
                              	(0,0) -- (31.31:0.6) arc (31.31:153.82:0.6) -- cycle;
                              \draw [shift={(-1,-1)},color=loko2!80!black,fill=loko2!80!black,fill opacity=0.1] 
                              	(0,0) -- (-33.69:0.6) arc (-33.69:88.82:0.6) -- cycle;
                              \draw [shift={(4.11,2.11)},fill=black,fill opacity=0.1] 
                              	(0,0) -- (-169.96:0.6) arc (-169.96:-148.69:0.6) -- cycle;
                              \draw [shift={(2,-3)},fill=black,fill opacity=0.1] 
                              	(0,0) -- (125.04:0.6) arc (125.04:146.31:0.6) -- cycle;
                              \draw [shift={(-0.95,1.21)},color=loko1,fill=loko1,fill opacity=0.1] 
                              	(0,0) -- (-91.18:0.6) arc (-91.18:-54.96:0.6) -- cycle;
                              \draw [shift={(-0.95,1.21)},color=loko1,fill=loko1,fill opacity=0.1] 
                              	(0,0) -- (-26.18:0.6) arc (-26.18:10.04:0.6) -- cycle;
                              \draw [color=loko1] (-1,-1)-- (2,-3);
                              \draw [color=loko1] (1.03,0.24)-- (4.11,2.11);
                              \draw [dash pattern=on 2pt off 2pt] (1.03,0.24)-- (-0.95,1.21);
                              \draw [dotted,color=loko2!80!black] (4.11,2.11)-- (-0.95,1.21);
                              \draw [dash pattern=on 2pt off 2pt] (-1,-1)-- (-0.95,1.21);
                              \draw [dotted,color=loko2!80!black] (2,-3)-- (-0.95,1.21);
                              \draw [shift={(4.11,2.11)}] (-169.96:0.6) arc (-169.96:-148.69:0.6);
                              \draw [shift={(4.11,2.11)}] (-169.96:0.5) arc (-169.96:-148.69:0.5);
                              \draw [shift={(2,-3)}] (125.04:0.6) arc (125.04:146.31:0.6);
                              \draw [shift={(2,-3)}] (125.04:0.5) arc (125.04:146.31:0.5);
                              \draw [shift={(-0.95,1.21)},color=loko1] (-91.18:0.6) arc (-91.18:-54.96:0.6);
                              \draw[color=loko1] (-0.76,0.71) -- (-0.72,0.59);
                              \draw[color=loko1] (-0.83,0.69) -- (-0.8,0.57);
                              \draw [shift={(-0.95,1.21)},color=loko1] (-26.18:0.6) arc (-26.18:10.04:0.6);
                              \draw[color=loko1] (-0.42,1.17) -- (-0.3,1.16);
                              \draw[color=loko1] (-0.43,1.1) -- (-0.31,1.08);
                              \fill [color=black] (-0.95,1.21) circle (1.5pt);
                              \draw[color=black] (-1.2,1.52) node {$I$};
                              \fill [color=black] (-1,-1) circle (1.5pt);
                              \draw[color=black] (-1.34,-1.12) node {$M$};
                              \fill [color=black] (2,-3) circle (1.5pt);
                              \draw[color=black] (1.74,-3.25) node {$N$};
                              \fill [color=black] (1.03,0.24) circle (1.5pt);
                              \draw[color=black] (1.36,0.12) node {$M'$};
                              \fill [color=black] (4.11,2.11) circle (1.5pt);
                              \draw[color=black] (4.18,1.8) node {$N'$};
                              \draw[color=loko2!80!black] (1.,1.1) node {$\alpha$};
                              \draw[color=loko2!80!black] (-.1,-0.84) node {$\alpha$};
                              \draw[color=black] (3.2,1.7) node {$\beta$};
                              \draw[color=black] (1.25,-2.25) node {$\beta$};
                              \draw[color=loko1] (-.65,.35) node {$\gamma$};
                              \draw[color=loko1] (0,1) node {$\gamma$};
                              \end{tikzpicture}
                           \caption{$M'N'$ est l'image de $MN$ par une rotation de centre $I$.}
                        \end{figure}
                  \item[Reflexion :]
\index{Isometrie@Isom\'etrie!reflexion@r\'eflexion}
                        Consid\'erons la sym\'etrie orthogonale d'axe $(\mathscr{D})$.
                        \begin{figure}[H]
                        \centering
                              \begin{tikzpicture}[line cap=round,line join=round,>=triangle 45,x=1.0cm,y=1.0cm]
                                    \draw [domain=-3.02:2.3] plot(\x,{(--1--2*\x)/3});
                                    \draw (-2.02,0.67)-- (-1.07,2.47);
                                    \draw (-0.46,-1.66)-- (1.56,-1.48);
                                    \draw [color=loko2] (-1.07,2.47)-- (0.25,0.5);
                                    \draw [color=loko2] (-0.34,1.54) -- (-0.49,1.44);
                                    \draw [color=loko2] (0.25,0.5)-- (1.56,-1.48);
                                    \draw [color=loko2] (0.98,-0.44) -- (0.83,-0.54);
                                    \draw [color=loko1] (-2.02,0.67)-- (-1.24,-0.49);
                                    \draw [color=loko1] (-1.58,0.17) -- (-1.72,0.07);
                                    \draw [color=loko1] (-1.54,0.11) -- (-1.69,0.01);
                                    \draw [color=loko1] (-1.24,-0.49)-- (-0.46,-1.66);
                                    \draw [color=loko1] (-0.8,-1) -- (-0.95,-1.1);
                                    \draw [color=loko1] (-0.76,-1.06) -- (-0.91,-1.16);
                                    \fill [color=black] (-2.02,0.67) circle (1.5pt);
                                    \draw[color=black] (-2.18,0.96) node {$M$};
                                    \fill [color=black] (-1.07,2.47) circle (1.5pt);
                                    \draw[color=black] (-0.9,2.74) node {$N$};
                                    \fill [color=black] (-0.46,-1.66) circle (1.5pt);
                                    \draw[color=black] (-0.34,-1.98) node {$M'$};
                                    \fill [color=black] (1.56,-1.48) circle (1.5pt);
                                    \draw[color=black] (1.76,-1.22) node {$N'$};
                                    \fill [color=black] (-1.24,-0.49) circle (1.5pt);
                                    \draw[color=black] (-1.58,-0.42) node {$A$};
                                    \fill [color=black] (0.25,0.5) circle (1.5pt);
                                    \draw[color=black] (0.62,0.5) node {$B$};
                                    \draw (1.9,2.1) node {$(\mathscr{D})$};
                              \end{tikzpicture}
                           \caption{$M'N'$ est l'image de $MN$ par la sym\'etrie orthogonale d'axe $(\mathscr{D})$.}
                        \end{figure}
                        Les trap\`ezes $ABNM$ et $M'N'BA$ sont superposable.
                  \item[Homoth\'etie :]
\index{Isometrie@Isom\'etrie!homothetie (de rapport +-1)@homoth\'etie (de rapport $\pm 1$)}
                        Consid\'erons l'homoth\'etie de centre $I$ et de rapport $k$.
                        
                        \begin{figure}[H]
                        \centering
                              \begin{tikzpicture}[line cap=round,line join=round,>=triangle 45,x=1.0cm,y=1.0cm]
                              \clip(-0.54,-1) rectangle (4.7,2.68);
                              \draw (2.34,-0.28)-- (1.86,1.12);
                              \draw (4.1,-0.49)-- (3.26,1.96);
                              \draw [dash pattern=on 1pt off 1pt] (0,0)-- (3.26,1.96);
                              \draw [dash pattern=on 1pt off 1pt] (0,0)-- (4.1,-0.49);
                              \fill [color=black] (0,0) circle (1.5pt);
                              \draw[color=black] (-0.18,0.22) node {$I$};
                              \fill [color=black] (2.34,-0.28) circle (1.5pt);
                              \draw[color=black] (2.52,-0.02) node {$M$};
                              \fill [color=black] (1.86,1.12) circle (1.5pt);
                              \draw[color=black] (1.76,1.42) node {$N$};
                              \fill [color=black] (4.1,-0.49) circle (1.5pt);
                              \draw[color=black] (4.3,-0.24) node {$M'$};
                              \fill [color=black] (3.26,1.96) circle (1.5pt);
                              \draw[color=black] (3.46,2.22) node {$N'$};
                              \end{tikzpicture}
                           \caption{$M'N'$ est l'image de $MN$ par une homoth\'etie de centre $I$.}
                        \end{figure}
                        
                        Comme $\overrightarrow{IM'} = k \vec{IM}$ et $\overrightarrow{IN'} = k \overrightarrow{IN}$, 
                        on a $\frac{IM'}{IM} = k$ et $\frac{IN'}{IN} = k$.
                        
                        Comme $\frac{IM'}{IM} = \frac{IN'}{IN}$, 
                        d'apr\`es la r\'eciproque de la propri\'et\'e de Thal\`es,
                        $(MN) \parallel (M'N')$.
                        On a alors $\frac{IM'}{IM} = \frac{IN'}{IN} = \frac{M'N'}{MN} = k$.
                        
                        Une homoth\'etie de rapport $k$ est une isom\'etrie si et seulement si $|k| = 1$.
            \end{description}
      \end{ex}
      %
      \begin{prop}[Propri\'et\'es des isom\'etries]
\index{Isometrie@Isom\'etrie!proprietes@propri\'et\'es}
            \begin{itemize}
                  \item La composition de deux isom\'etries est une isom\'etrie.
                  \item La r\'eciproque d'une isom\'etrie est une isom\'etrie.
                  \item \begin{itemize}
                              \item L'image d'une droite par une isom\'etrie est une droite.
                              \item Toute isom\'etrie conserve l'alignement.
                        \end{itemize}
                  \item Toute isom\'etrie transforme deux droites parall\`eles en deux droites parall\`eles :
                        on dit qu'une isom\'etrie \emph{conserve le parall\'elisme}.
                  \item Toute isom\'etrie transforme deux droites perpendiculaires en deux droites perpendiculaires :
                        on dit qu'une isom\'etrie \emph{conserve l'orthogonalit\'e}.
                  \item Toute isom\'etrie transforme un cercle en un cercle : 
                        on dit qu'une isom\'etrie \emph{conserve la cocyclicit\'e}.
                  \item Toute isom\'etrie transforme 2 cercles tangents 
                        (respectivement un cercle et une droite tangente)
                        en 2 cercles tangents
                        (respectivement en un cercle et une droite tangente) :
                        on dit qu'une isom\'etrie \emph{conserve le contact}.
            \end{itemize}
      \end{prop}
      %
\section{Reconnaissance des isom\'etries}
      \subsection{\`A l'aide de l'expression analytique}
            Soit $f$ une application affine d'expression analytique 
            \[
                  \left\lbrace
                        \begin{array}{l}
                              x' = a_1 x + a_2 y + a_3 \\
                              y' = b_1 x + b_2 y + b_3 \;.
                        \end{array}
                  \right.
            \]
            Soit 
            $ A = \begin{pmatrix} a_1 & a_2 \\ b_1 & b_2 \end{pmatrix} $
            et soit $\Delta = a_1 \cdot b_2 - b_1 \cdot a_2$
            .
            %
            \begin{prop}
                  L'application $f$ est une isom\'etrie si et seulement si $A$ est de l'un des types : 
                  \begin{align*}
                        \begin{pmatrix}
                              a & -b \\ b & a 
                        \end{pmatrix}
                              \qquad \text{ ou } \qquad
                        \begin{pmatrix}
                              a & b \\ b & -a
                        \end{pmatrix}
                  \end{align*}
                  o\`u $a^2 + b^2 = 1$.
                  
                  De plus, $f$ est une isom\'etrie positive (ou un d\'eplacement) si $\Delta = 1$
                  et elle est une isom\'etrie n\'egative (ou un antid\'eplacement) si $\Delta = -1$.
            \end{prop}
      \subsection{A l'aide de l'expression complexe}
            Soit $f$ une application affine.
            \begin{itemize}
                  \item L'application $f$ est une isom\'etrie positive si son expression complexe est de la forme
                        $z' = a z + b$ o\`u $|a| = 1$.
                  \item L'application $f$ est une isom\'etrie n\'egative si son expression complexe est de la forme
                        $z' = a \bar{z} + b$ o\`u $|a| = 1$.
            \end{itemize}
      
      \subsection{Selon le nombre de points invariants}
      %
            Soit $f$ une isom\'etrie.
		      \begin{table}[h!]
			            \begin{tabularx}{\linewidth}{|p{.27\linewidth}|p{.18\linewidth}|X|}
                        \hline
                              \textbf{Points invariants} 
                                    & \textbf{d\'eplacement}   
                                          & \textbf{antid\'eplacement}
                                                \\
                        \hline
                              Le plan $\Plan$ tout entier            
                                    & $\Id$                    
                                          & 
                                                \\
                        \hline
                              Une droite $\Droite$                    
                                    &                          
                                          & $S_{\Droite}$
                                                \\
                        \hline
                              Un point $\Omega$                       
                                    & $r(\Omega,\theta),\ \theta\neq 0$ 
                                          & 
                                                \\
                        \hline
                              Aucun point invariant                   
                                    & $t_{\vec{u}},\ \vec{u}\neq\vec{0}$ 
                                          & \raggedright
                                            Compos\'ee d'une sym\'etrie orthogonale d'axe $\Droite$ et 
                                            d'une translation de vecteur $\vec{u}$ non nul de 
                                            direction parall\`ele \`a $\Droite$.
                                                 \tabularnewline[0\baselineskip]
                        \hline
		            \end{tabularx}
            \caption{Classification des isom\'etries suivant l'ensemble des points invariants.}
            \label{tab:classification des isometries}
            \end{table}
%
%      \begin{thm}
%\index{Isometrie@Isom\'etrie!reconnaissance des}
%            Soit $f$ une isom\'etrie.
%            \begin{itemize}
%                  \item Si $f$ fixe trois points non-align\'es alors $f$ est l'application identit\'e, not\'ee $\Id$.
%                  \item Si $f \neq \Id$ et $f$ fixe deux points distincts  $A$ et $B$ 
%                        alors $f$ est la sym\'etrie d'axes $(AB)$.
%                  \item Si $f$ a un seul point invariant alors $O$ alors $f$ est la rotation de centre $O$ 
%                        et d'angle $\alpha$ ($\alpha \not = 0 \quad [2\pi]$).
%            \end{itemize}
%            %
%            \begin{rmk}
%                  Une translation est une isom\'etrie qui n'a aucun point invariant
%                  mais toutes isom\'etrie qui n'a aucun point invariant n'est pas forc\'ement une translation.
%            \end{rmk}
%      \end{thm}
      %
\section{Classification des isom\'etries}
      Soit $\widehat{OAB}$ un angle orient\'e.
      \begin{itemize}
            \item La translation ne change pas l'orientation des angles.
                  \begin{figure}[H]
                  \centering
                        \begin{tikzpicture}[line cap=round,line join=round,>=stealth,x=1.0cm,y=1.0cm]
                              \draw [shift={(-2.82,-1.26)}] 
                                    (0,0) -- (-8.25:0.49) arc (-8.25:39.74:0.49) -- cycle;
                              \draw [shift={(1.05,-1.25)}] 
                                    (0,0) -- (-8.25:0.49) arc (-8.25:39.74:0.49) -- cycle;
                              \draw [shift={(-2.82,-1.26)},->,color=loko1] (-8.25:0.49) arc (-8.25:39.74:0.49);
                              \draw[color=loko1] (-0.54,-1.59)-- (-2.82,-1.26);
                              \draw[color=loko1] (-2.82,-1.26)-- (-1.35,-0.04);
                              \draw [->] (-1.41,0.81) -- (2.46,0.82);
                              \draw[color=loko2] (3.32,-1.58)-- (1.05,-1.25);
                              \draw[color=loko2] (1.05,-1.25)-- (2.51,-0.03);
                              \draw[color=loko2] [shift={(1.05,-1.25)},->] (-8.25:0.49) arc (-8.25:39.74:0.49);
                              \begin{scriptsize}
                              \fill [color=loko1] (-0.54,-1.59) circle (1.5pt) node[above] {$A$};
                              \fill [color=loko1] (-2.82,-1.26) circle (1.5pt) node[above] {$O$};
                              \fill [color=loko1] (-1.35,-0.04) circle (1.5pt) node[above] {$B$};
                              \draw[color=black] (1.04,1) node {$\vec{u}$};
                              \fill [color=loko2] (3.32,-1.58) circle (1.5pt) node[below] {$A'$};
                              \fill [color=loko2] (1.05,-1.25) circle (1.5pt) node[below] {$O'$};
                              \fill [color=loko2] (2.51,-0.03) circle (1.5pt) node[below] {$B'$};
                              \end{scriptsize}
                        \end{tikzpicture}
                     \caption{La translation est une \emph{isom\'etrie positive}.}
                  \end{figure}
            \item La rotation ne change pas l'orientation des angles.
                  \begin{figure}[H]
                  \centering
                        \begin{tikzpicture}[line cap=round,line join=round,>=stealth,x=1.0cm,y=1.0cm]
%                              \draw [shift={(2.36,-0.93)},color=loko2!50!orange] 
%                                    (0,0) -- (131.26:0.49) arc (131.26:160.43:0.49) -- cycle;
%                              \draw [shift={(2.2,-1.7)},color=loko1] 
%                                    (0,0) -- (176.26:0.49) arc (176.26:205.43:0.49) -- cycle;
                              \draw [shift={(3.22,-1.51)},double,->,line width=0.4pt,color=loko3!50!black] 
                                    (146.06:0.52) arc (146.06:191.06:0.52); % -- cycle;
                              \draw [shift={(2.36,-0.93)},->,color=loko2] (131.26:0.49) arc (131.26:160.43:0.49);
                              \draw [color=loko2] (0.71,0.95)-- (2.36,-0.93);
                              \draw [color=loko2] (2.36,-0.93)-- (0.88,-0.4);
                              \draw [color=loko1] (-0.29,-1.54)-- (2.2,-1.7);
                              \draw [color=loko1] (2.2,-1.7)-- (0.78,-2.38);
                              \draw [shift={(2.2,-1.7)},->,color=loko1] (176.26:0.49) arc (176.26:205.43:0.49);
                              \draw [shift={(3.22,-1.51)},line width=0.2pt,dashed]  
                                    plot[domain=2.37:3.15,variable=\t]
                                          ({1*3.51*cos(\t r)+0*3.51*sin(\t r)},{0*3.51*cos(\t r)+1*3.51*sin(\t r)});
                              \draw [shift={(3.22,-1.51)},line width=0.2pt,dashed] 
                                    plot[domain=2.7:3.49,variable=\t]
                                          ({1*2.59*cos(\t r)+0*2.59*sin(\t r)},{0*2.59*cos(\t r)+1*2.59*sin(\t r)});
                              \draw [shift={(3.22,-1.51)},line width=0.2pt,dashed]  
                                    plot[domain=2.55:3.33,variable=\t]
                                          ({1*1.03*cos(\t r)+0*1.03*sin(\t r)},{0*1.03*cos(\t r)+1*1.03*sin(\t r)});
                              \draw [line width=0.2pt,dotted] (3.22,-1.51)-- (2.36,-0.93);
                              \draw [line width=0.2pt,dotted] (2.2,-1.7)-- (3.22,-1.51);
                              \begin{scriptsize}
                                    \fill [color=loko2] (0.71,0.95) circle (1.5pt);
                                    \draw[color=loko2] (0.84,1.16) node {$A$};
                                    \fill [color=loko2] (2.36,-0.93) circle (1.5pt);
                                    \draw[color=loko2] (2.49,-0.72) node {$O$};
                                    \fill [color=loko2] (0.88,-0.4) circle (1.5pt);
                                    \draw[color=loko2] (1.01,-0.19) node {$B$};
                                    \fill [color=loko2] (3.22,-1.51) circle (1.5pt);
                                    \draw[color=loko2] (3.3,-1.29) node {$I$};
                                    \fill [color=loko1] (-0.29,-1.54) circle (1.5pt);
                                    \draw[color=loko1] (-0.15,-1.33) node {$A'$};
                                    \fill [color=loko1] (2.2,-1.7) circle (1.5pt);
                                    \draw[color=loko1] (2.09,-1.94) node {$O'$};
                                    \fill [color=loko1] (0.78,-2.38) circle (1.5pt);
                                    \draw[color=loko1] (0.92,-2.17) node {$B'$};
                              \end{scriptsize}
                        \end{tikzpicture}
                     \caption{La rotation est une \emph{isom\'etrie positive}.}
                  \end{figure}
            \item La r\'eflexion change l'orientation des angles.
                  \begin{figure}[H]
                  \centering
                        \begin{tikzpicture}[line cap=round,line join=round,>=stealth,x=1.0cm,y=1.0cm]
                        \clip(-2.72,-3.6) rectangle (2.1,.5);
                              \draw [shift={(-2.37,-2.1)},color=loko2,fill=loko2,fill opacity=0.1] 
                                    (0,0) -- (17.18:0.49) arc (17.18:47.15:0.49) -- cycle;
                              \draw [shift={(-1.59,-3.28)},color=loko1,fill=loko1,fill opacity=0.1] 
                                    (0,0) -- (19.91:0.49) arc (19.91:49.89:0.49) -- cycle;
                              \draw [shift={(-2.37,-2.1)},->,color=loko2] (17.18:0.49) arc (17.18:47.15:0.49);
                              \draw [color=loko2] (-0.56,-1.54)-- (-2.37,-2.1);
                              \draw [color=loko2] (-2.37,-2.1)-- (-0.26,0.17);
                              \draw [domain=-2.72:1.91] plot(\x,{(--3.78-1.81*\x)/-2.74});
                              \draw [color=loko1] (-0.36,-1.83)-- (-1.59,-3.28);
                              \draw [color=loko1] (-1.59,-3.28)-- (1.33,-2.23);
                              \draw [shift={(-1.59,-3.28)},<-,color=loko1] (19.91:0.49) arc (19.91:49.89:0.49);
                              \begin{scriptsize}
                                    \fill [color=loko2] (-0.56,-1.54) circle (1.5pt);
                                    \draw[color=loko2] (-0.43,-1.33) node {$A$};
                                    \fill [color=loko2] (-2.37,-2.1) circle (1.5pt);
                                    \draw[color=loko2] (-2.09,-2.24) node {$O$};
                                    \fill [color=loko2] (-0.26,0.17) circle (1.5pt);
                                    \draw[color=loko2] (-0.13,0.39) node {$B$};
                                    \fill [color=loko1] (-0.36,-1.83) circle (1.5pt);
                                    \draw[color=loko1] (-0.16,-1.72) node {$A'$};
                                    \fill [color=loko1] (-1.59,-3.28) circle (1.5pt);
                                    \draw[color=loko1] (-1.28,-3.45) node {$O'$};
                                    \fill [color=loko1] (1.33,-2.23) circle (1.5pt);
                                    \draw[color=loko1] (1.49,-2.02) node {$B'$};
                                    \draw[color=black] (1.7,.2) node {$(\mathscr{D})$};
                              \end{scriptsize}
                        \end{tikzpicture}
                     \caption{La r\'eflexion est une \emph{isom\'etrie n\'egative}.}
                  \end{figure}
      \end{itemize}
      %
      Toute isom\'etrie qui \emph{conserve l'orientation des angles} est appel\'ee \emph{d\'eplacement}.
\index{D\'eplacement}      
      Toute isom\'etrie qui \emph{\textbf{ne} conserve \textbf{pas} l'orientation des angles} est appel\'ee
\index{Antid\'eplacement}      
       \emph{\textbf{anti}d\'eplacement}.
      %
      \begin{prop}
            \begin{itemize}
                  \item Toute isom\'etrie est soit un d\'eplacement, soit un antid\'eplacement.
                  \item La r\'eciproque d'un d\'eplacement est un d\'eplacement.
                  \item La r\'eciproque d'un antid\'eplacement est un antid\'eplacement.
                  \item La composition de deux d\'eplacements est un d\'eplacement.
                  \item La composition de deux antid\'eplacements est un d\'eplacement.
                  \item La composition d'un d\'eplacement et d'un antid\'eplacement,
                        dans n'importe quel ordre, est un antid\'eplacement.
            \end{itemize}
      \end{prop}
      %
      \begin{prop}[Expression complexe d'un d\'eplacement]
\index{D\'eplacement!Expression complexe}
	      Tout d\'eplacement a pour expression complexe une \'equation de la forme $z' = az + b$ avec $|a| = 1$.
	            \begin{center}
	            \resizebox{0.96\linewidth}{!}{$\xymatrix@C=.05\actualtextwidth{
                            &
                                   &  
                                          & *+[F]{ \begin{array}{c} f:z' = az + b \\ |a| = 1 \end{array} }
                                            \ar[dr]
                                            \ar[dll] 
                                                 &
                                                        \\
                            & *+[F]{a=1} 
                              \ar[dr]
                              \ar[dl]
                                   &  
                                          &
                                                 & *+[F]{a \neq 1 }
                                                   \ar[dd]
                                                        \\
                     *+[F]{b=0} 
                     \ar[d]
                            &            
                                   & *+[F]{b\neq 0}
                                     \ar[d]
                                          &
                                                 &
                                                        \\ 
                     *+[F]{f=\Id}
                            & 
                                   & *+[F]{\begin{array}{l} f=t_{\vec{u}} \\ \vec{u} (b) \end{array}}
                                          & 
                                                 & *+[F]{\begin{array}{l} f=r(I;\alpha) \\ \alpha= \Arg{a} \\ z_I=\frac{b}{1-a} \end{array} } \\
	                     }$}
	            \end{center}
      \end{prop}
      %
      \begin{prop}[Expression complexe d'un antid\'eplacement]
\index{Antid\'eplacement!Expression complexe}
	            Tout antid\'eplacement a pour expression complexe une \'equation de la forme :
                  $z' = a\bar{z} + b$.
            \nota{
                  %\begin{itemize}
                        %\item 
                              Si $a\bar{b}+b = 0$, $f$ est une r\'eflexion. \newline
                        %\item 
                              Si $a\bar{b}+b \neq 0$, il existe une droite $(\mathscr{D})$ et un vecteur 
	                              $\vec{u} \neq \vec{0}$ qui dirige $(\mathscr{D})$ tels que 
                              $f = t_{\vec{u}} \circ S_{(\mathscr{D})}$
                              (forme r\'eduite d'un antid\'eplacement).
                  %\end{itemize}
            }
      \end{prop}
      
%
\section{Composition de d\'eplacements et antid\'eplacements}
\index{Composition!deux antideplacements@deux antid\'eplacements}
      On d\'esigne par :
      \begin{description}
            \item[$s_{\Droite}$] la sym\'etrie orthogonal d'axe $\Droite$,
            \item[$t_{\vec{u}}$] la translation de vecteur $\vec{u}$,
            \item[$r(\Omega,\theta)$] la rotation de centre $\Omega$ et d'angle $\theta$.
      \end{description}
      %
      \subsection{Composition de deux sym\'etrie d'axes parall\`eles}
\index{Translation!composition de deux symetries@composition de deux sym\'etries}
            Soient $\DroiteU$ et $\DroiteD$ deux droites parall\`eles distinctes.
            Soient $H_1$ un point de $\DroiteU$ et $H_2$ le projet\'e orthogonal de $H_1$ sur $\DroiteD$.
            Alors \thmbox{$s_{\DroiteD} \circ s_{\DroiteU} = t_{2\overrightarrow{H_1 H_2}}$}.
            
            La r\'eciproque est aussi vraie.
            Soit $\vec{u}$ un vecteur et $\DroiteU$ une droite
            perpendiculaire \`a la direction de $\vec{u}$. 
            Soit $\DroiteD$ le translat\'e de $\DroiteU$ par le vecteur $\frac{1}{2}\vec{u}$.
            Alors \thmbox{$t_{\vec{u}} = s_{\DroiteD} \circ s_{\DroiteU}$}.
            \begin{figure}[h!]
            		\centering
                  \begin{tikzpicture}[line cap=round,line join=round,>=triangle 45,x=1.0cm,y=1.0cm]
                  \clip(-1.96,-0.1) rectangle (6.1,5.28);
                  \draw [->] (-1.48,3.76) -- (1.84,2.82);
                  \draw [domain=-1.96:6.1] plot(\x,{(-5.9--3.32*\x)/0.94});
                  \draw [domain=-1.96:6.1] plot(\x,{(-11.85--3.32*\x)/0.94});
                  \begin{scriptsize}
                  \draw[color=black] (0.82,3.5) node {$\vec{u}$};
                  \draw[color=black] (4.45,5) node {$\DroiteD$};
                  \draw[color=black] (2.7,5) node {$\DroiteU$};
                  \fill [color=black] (-0.68,2.3) circle (1.5pt);
                  \draw[color=black] (-0.52,2.56) node {$A$};
                  \fill [color=black] (5.07,0.67) circle (1.5pt);
                  \draw[color=black] (5.26,0.94) node {$A'$};
                  \fill [color=black] (2.64,1.36) circle (1.5pt);
                  \draw[color=black] (2.86,1.62) node {$A''$};
                  \end{scriptsize}
                  \end{tikzpicture}
               \caption{La compos\'ee de deux sym\'etries d'axes parall\`eles est une translation.}
            \label{fig:translation = compose de deux symetries}
            \end{figure}
      %
      \subsection{Composition de deux sym\'etrie d'axes s\'ecantes}
\index{Rotation!composition de deux symetries@composition de deux sym\'etries}
            Soient $\DroiteU$ et $\DroiteD$ deux droites distinctes se coupant en un point $\Omega$.
            Soit $\theta = \mathrm{mes}(\DroiteU,\DroiteD) \quad [\pi]$.
            Alors \thmbox{$s_{\DroiteD} \circ s_{\DroiteU} = r(\Omega,2\theta)$}.
            
            La r\'eciproque est aussi vraie.
            Soit $\Omega$ un point et $\theta$ un nombre r\'eel\footnote{qui sera consid\'er\'e en tant qu'angle}.
            Soit $\DroiteU$ une droite et soit 
            $\DroiteD$ le transform\'e de $\DroiteU$ par la rotation de centre $\Omega$ et d'angle $\frac{\theta}{2}$.
            Alors \thmbox{$r(\Omega,\theta) = s_{\DroiteD} \circ s_{\DroiteU}$}.
            
            \begin{figure}[h!]
            \centering
                  \begin{tikzpicture}[line cap=round,line join=round,>=triangle 45,x=1.0cm,y=1.0cm]
                  \clip(3.2,0.2) rectangle (14.66,8);
                  \draw(9.64,5.08) -- (9.49,5.32) -- (9.25,5.16) -- (9.41,4.93) -- cycle; 
                  \draw(10.87,1.99) -- (10.59,2) -- (10.58,1.71) -- (10.86,1.71) -- cycle; 
                  \draw [shift={(4.78,1.86)},fill=black,fill opacity=0.1] (0,0) -- (-10.65:0.6) arc (-10.65:59.35:0.6) -- cycle;
                  \draw [shift={(4.78,1.86)}] (0,0) -- (178.55:0.4) arc (178.55:213.55:0.4) -- cycle;
                  \draw [domain=3.2:14.66,loko2] plot(\x,{(-7.84--0.1*\x)/-3.96});
                  \draw [domain=3.2:14.66,loko1] plot(\x,{(-4.33--2.19*\x)/3.3});
                  \draw [shift={(4.78,1.86)}] 
                        plot[domain=-0.19:1.04,variable=\t]
                              ({1*6.17*cos(\t r)+0*6.17*sin(\t r)},{0*6.17*cos(\t r)+1*6.17*sin(\t r)});
                  \draw[loko2] (14,2) node {$ \DroiteU $};
                  \draw[loko1] (14,7.5) node {$ \DroiteD $};
                  \draw [dash pattern=on 3pt off 3pt] (10.89,2.69)-- (10.84,0.72);
                  \draw [dash pattern=on 3pt off 3pt] (10.89,2.69)-- (7.92,7.16);
                  \draw [shift={(4.78,1.86)}] (-10.65:0.6) arc (-10.65:59.35:0.6);
                  \draw [shift={(4.78,1.86)}] (-10.65:0.5) arc (-10.65:59.35:0.5);
                  \draw[dotted] (7.92,7.16)-- (4.78,1.86);
                  \draw[dotted] (4.78,1.86)-- (10.84,0.72);
                  \begin{scriptsize}
                  \fill [color=black] (4.78,1.86) circle (1.5pt) node[below] {$\Omega$};
                  \fill [color=black] (10.84,0.72) circle (1.5pt) node[right] {$A$};
                  \fill [color=black] (10.89,2.69) circle (1.5pt) node[right] {$A'$};
                  \fill [color=black] (7.92,7.16) circle (1.5pt) node[above right] {$A''$};
                  \end{scriptsize}
                  \end{tikzpicture}
               \caption{La compos\'ee de deux sym\'etries d'axes s\'ecantes est une rotation.}
            \end{figure}
      %
      \subsection{Composition de deux d\'eplacements}
\index{Composition!de deux rotations}
            \paragraph{Compos\'ee de deux rotations}
                  Soient $r_1 = r(\Omega , \theta_1)$ et $r_2 = r(\Omega_2 , \theta_2)$.
                  \begin{itemize}
                        \item Si $\theta_1 + \theta_2 = 0 \quad [2\pi]$ alors $r_2 \circ r_1$ est une translation.
                        \item Si $\theta_1+\theta_2\neq 0\quad [2\pi]$ alors $r_2\circ r_1$ est 
                              une rotation d'angle $\theta_1+\theta_2$.
                  \end{itemize}
            \paragraph{Compos\'ee d'une rotation et d'une translation}
                  Soient $t$ une translation et $r$ une rotation d'angle $\theta \neq 0 \quad [2\pi]$.
                  Alors la $r \circ t$ est une rotation d'angle $\theta$.
      %
      \subsection{Conclusion}
            \begin{itemize}
                  \item Tout d\'eplacement est la compos\'ee de deux sym\'etries orthogonales.
                  \item Tout anti-d\'eplacement est soit une sym\'etrie orthogonale, 
                        soit la compos\'ee de trois sym\'etries orthogonales.
            \end{itemize}
            %
            \smallnota{Une sym\'etrie centrale de centre $\Omega$ est la m\^eme chose qu'une
                        rotation de centre $\Omega$ et d'angle $\pi$.}
            %
















\TCDLessonStyleOff
\clearpage
% Isometries %
\providecommand{\Droite}{(\mathscr{D})}
\providecommand{\AngleEgal}[3]{#1 = #2 \quad \left[ #3 \right] }
\providecommand{\mes}[2]{\mathrm{mes}\left( \widehat{ \overrightarrow{#1} , \overrightarrow{#2} } \right)}
\section{Exercices}
























